\paragraph{About this part.} Part~II asked what the book predicts. Part~III asks what a trader can do
with a prediction once the mechanics of trading are included. It compares crossing the spread with
waiting in the queue (Section~\ref{sec:passive}); shows that an order's place in the queue determines
whether, and when, it is filled (Section~\ref{sec:queue-position}); explains why the fills a passive
order does get are biased towards bad ones (Section~\ref{sec:adverse}); measures how latency eats into a
signal (Section~\ref{sec:latency}); and states what a simulation must do to answer these questions
honestly (Section~\ref{sec:replay}).

% ================================================================
\section{Passive and aggressive execution}
\label{sec:passive}
% ================================================================

Consider one unit traded on side $d$ at time $t$. For an aggressive order, which executes
immediately, a simplified realised profit is
\begin{equation}
  \Pi^{\text{agg}} = d\,\Delta m_{t,H} - C_{\text{spread}} - C_{\text{fees}} - C_{\text{impact}},
\end{equation}
where $C_{\text{spread}} = s_t/2$ is the half spread paid to cross, $C_{\text{fees}}$ the exchange
and clearing fees, and $C_{\text{impact}}$ the additional price concession from walking the book.
\inwords{profit from trading immediately is the price move in our favour---multiplying by $d$
flips the sign for a sell, so a falling price is good for a seller---minus what it cost to get in:
half the spread, the fees, and any extra paid for consuming more than the best level.}
A passive order placed at $t$ trades only if it is filled, at time $t_F$ and price
$p^{\text{fill}}$. Define its post-fill mark-out over horizon $h$ as
\begin{equation}
  R^{\text{fill}}_{h} = d\left(m_{t_F+h} - p^{\text{fill}}\right).
  \label{eq:markout}
\end{equation}
\inwords{once our order has traded, look at the mid $h$ seconds later and ask how far it is from
the price we got, counted positive when it is in our favour. A large positive mark-out means we
bought cheaply; a negative one means the market moved against us right after we traded.}
For a passive buy filled at the bid, the mark-out is the half spread at the fill time,
$s_{t_F}/2$, plus the mid move after the fill. The passive profit is
\begin{equation}
  \Pi^{\text{pass}} = \mathbb{1}_{\Fill}\left(R^{\text{fill}}_{h} + R_{\text{rebate}}\right),
\end{equation}
where $R_{\text{rebate}}$ is any liquidity-provision rebate net of fees.
\inwords{if the order is filled, the profit is its mark-out plus any rebate the exchange pays for
providing liquidity; if it is never filled, the indicator is zero and so is the profit.} The passive profit exists
only conditional on execution, which gives the decomposition
\begin{equation}
  \E[\Pi^{\text{pass}}] = \P(\Fill)\; \E[\Pi^{\text{pass}} \mid \Fill].
  \label{eq:decomp}
\end{equation}
\inwords{the average profit of a passive order equals the chance it is filled times the average
profit when it is. This is the identity $\E[Z\,\mathbb{1}_A] = \P(A)\,\E[Z \mid A]$ of
Section~\ref{sec:primer}, with $A$ the fill.}
Equation~\eqref{eq:decomp} separates two quantities that a directional forecast conflates: how
likely the order is to trade, and how good the trade is when it does. The rest of this part is
about the two factors. Market-making models~\citep{avellaneda2008,cartea2015} and optimal-execution
models~\citep{almgren2001} give the surrounding control problems; here we focus on the inputs
those problems need from the book.

\subsection{Signal value and execution value}
\label{sec:tradeability-framework}

The decomposition gives tradeability a precise meaning. A directional model estimates the
\emph{signal value}
\begin{equation}
  V^{\text{sig}}(X_t) = \E[\Delta m_{t,H} \mid X_t],
\end{equation}
the expected price move given what is known. A passive execution decision is worth its
\emph{execution value}
\begin{equation}
  V^{\text{exec}}(X_t, a_t) = \P(\Fill \mid X_t, a_t)\; \E\big[\Pi^{\text{pass}} \mid \Fill, X_t, a_t\big],
\end{equation}
which depends on the action through where the order is placed and therefore on its queue position.
\inwords{signal value asks ``which way will the price go?''; execution value asks ``if I act on that
view with this order, how likely am I to trade, and how good will the trade be when I do?'' A model
can be excellent at the first and useless at the second.}

The two can even point in opposite directions. A forecast of a falling price is most confident when
the book is already one-sided---a thin bid, heavy selling. That is exactly when a resting buy order is
most likely to be filled and least likely to be worth having: the fill and the adverse move arrive
together (Section~\ref{sec:queue-models}). Conditioning on the fill reverses the apparent value of the
signal. This is why the survey treats $\P(\Fill \mid \cdot)$ and $\E[\,\cdot \mid \Fill]$ as separate
modelling targets, and why evaluation must use the conditional quantities rather than unconditional
accuracy.

Who earns the execution value in practice is documented. A large high-frequency market maker studied
by \citet{menkveld2013} traded passively in about four of every five trades, earning the spread while
losing on its inventory; across high-frequency firms, differences in relative speed explain large
differences in trading performance~\citep{baron2019}.

% ================================================================
\section{Queue position and fill probability}
\label{sec:queue-position}
% ================================================================

A passive order is filled only when everything ahead of it in its queue has been traded or
cancelled. This section shows why that makes queue position a variable in its own right: two orders
at the same price, with the same forecast behind them, can have very different chances of being
filled.

\subsection{Residual queue position}

For a resting order, the displayed size at its price does not determine its fill probability. Let
$Q^{\text{ahead}}_t$ be the volume ahead of the order at the same price. Under price--time
priority the order cannot execute until that volume has been removed by executions or
cancellations. Two orders at the same price can therefore have very different values, and queue
position has been modelled explicitly as a source of value~\citep{moallemi2017}.

For a passive order with $q$ ahead of it, the probability of a fill within horizon $H$ is
$\P(\Fill \mid Q^{\text{ahead}}_t = q,\, X_t)$. It must be distinguished from
$\P(\Delta m_{t,H} > 0 \mid X_t)$. A model can predict direction excellently and still have poor
trading value if the predicted moves happen mainly when the passive order is unlikely to execute.

\begin{example}[Same forecast, different orders]
\label{ex:queue}
Two resting bids at the same price have 20 and 480 lots ahead of them. Net depletion runs at 40
lots per second. The naive time to reach the front is $0.5$ s for the first and $12$ s for the
second. Both carry the same $\P(\Delta m_{t,H} > 0 \mid X_t)$. Over a five-second horizon the first
almost surely fills and the second almost surely does not.
\end{example}

% ================================================================
\section{Adverse selection}
\label{sec:adverse}
% ================================================================

Being filled is not always good news. A passive order tends to be filled exactly when someone better
informed wants to trade against it, so the fills it gets are, on average, followed by price moves
against it. This section defines that cost precisely and shows how it depends on queue position.

\subsection{Conditioning on the fill}

Using the mark-out $R^{\text{fill}}_{h}$ of~\eqref{eq:markout}, define the conditional
adverse-selection component for an order that had $q$ ahead of it when placed as
\begin{equation}
  \mathrm{AS}(q, h) = -\E\left[R^{\text{fill}}_{h} - \tfrac{s_{t_F}}{2} \,\middle|\, \Fill,\, Q^{\text{ahead}}_t = q\right],
\end{equation}
the expected mid move against the order after it is filled, excluding the half spread it earns.
\inwords{among orders that started with $q$ lots ahead of them \emph{and were filled}, how far,
on average, did the price then move against us? It is a form of the winner's curse: in an auction,
the times your low bid wins are disproportionately the times when others knew something you did
not. A passive order is most likely to be filled just when an informed trader wants to trade the
other way.}
The conditioning on the fill is essential. Without it, the quantity measures the unconditional
drift of the market rather than the information in the event that the order actually traded.
Classical models make the mechanism explicit: a liquidity provider who trades with a better-informed
counterparty loses on average~\citep{glosten1985,kyle1985}.

\subsection{Queue position as an information variable}

Queue position affects both the probability of execution and the information content of
execution. \citet{moallemi2017} model both effects and find that a position at the front of the
queue is always worth more than one at the back: in their model adverse selection increases with queue
position, because an order at the end of a long queue tends to be filled only by a large aggressive
order, which is more likely to be informed. For some large-tick US stocks they estimate the front--back
value difference to be comparable to the bid--ask spread. The direction of the effect is therefore
supported by a model; its size is an empirical object that can depend on asset, tick regime,
spread, order-flow state, horizon and cancellation behaviour, and execution-aware benchmarks should
report it directly (Section~\ref{sec:agenda}).

\begin{example}[The decomposition at two positions]
\label{ex:as}
All values in ticks; fill window $H = 10$ s, mark-out horizon $h = 10$ s, one-tick spread, no
rebate. Here ``$\mid q$'' abbreviates $\mid Q^{\text{ahead}}_t = q$, and

\[
  \E[\Pi^{\text{pass}} \mid \Fill, q] = s_{t_F}/2 - \mathrm{AS}(q,h):
\]

\begin{center}\small
\begin{tabular}{lccccc}
\toprule
Position & $\P(\Fill \mid q)$ & $s_{t_F}/2$ & $-\mathrm{AS}(q,h)$ & $\E[\Pi^{\text{pass}} \mid \Fill, q]$ & $\E[\Pi^{\text{pass}} \mid q]$ \\
\midrule
Front & 0.80 & $+0.5$ & $-0.20$ & $+0.30$ & $+0.24$ \\
Back  & 0.25 & $+0.5$ & $-0.45$ & $+0.05$ & $+0.0125$ \\
\bottomrule
\end{tabular}
\end{center}
The table shows the arithmetic of~\eqref{eq:decomp}; the ordering of the two rows matches the
direction found by~\citet{moallemi2017}, with stylised magnitudes. A metric that looked
only at direction or only at fill rate would rank the two positions without seeing the product.
\end{example}

% ================================================================
\section{Latency and signal decay}
\label{sec:latency}
% ================================================================

A prediction at horizon $H$ is useful only if its information survives until the decision is
executed. Write the total latency as
\begin{equation}
  T_{\text{lat}} = T_{\text{signal}} + T_{\text{decision}} + T_{\text{network}} + T_{\text{exchange}},
\end{equation}
the times taken to compute the signal, to decide on an action, to transmit the order, and for the
exchange to process it.
\inwords{total delay from seeing the data to the order being live at the exchange is the sum of
the stage delays, just as the latency of a processor pipeline is the sum of its stages.}
A signal whose value decays materially before $T_{\text{lat}}$ is not economically equivalent to
one of the same accuracy with longer persistence. The relevant object is not
$\E[\Delta m_{t,H} \mid X_t]$ but
\begin{equation}
  \E\left[\Pi \mid X_t,\, T_{\text{lat}},\, \text{execution policy}\right],
\end{equation}
the expected profit given the information, the latency, and the rule that maps forecasts to orders.
\inwords{the quantity that matters is not ``will the price go up?'' but ``how much do we expect to
make, given what we knew, how slow we are, and what we do with the forecast?''}
Competition on latency is itself an economic phenomenon~\citep{budish2015,hasbrouck2013}.

\begin{example}[Half-life against latency]
Let a signal's value decay exponentially with a half-life of 2 ms, so a fraction
$2^{-T_{\text{lat}}/2\,\mathrm{ms}}$ of it remains after a delay $T_{\text{lat}}$. At $100\,\mu$s total latency
$2^{-0.05} \approx 97\%$ of its value remains; at 5 ms only $2^{-2.5} \approx 18\%$ does. The same
forecast accuracy yields very different economics.
\end{example}

% ================================================================
\section{Replay and deterministic simulation}
\label{sec:replay}
% ================================================================

The questions of this part---will the order fill, and is the fill a good one---cannot be answered by
a backtest that assumes orders fill when the price touches them. This section shows how such a
backtest misleads and lists what a simulation must model to give trustworthy answers; Part~IV
describes a system built to those requirements.

\subsection{Assumptions of a naive backtest}

A common backtest of a passive signal assumes that a limit order placed at the bid is filled when
the signal fires, or when the price touches the order's level:
\begin{lstlisting}
naive:   signal at t -> assume filled at bid -> P&L = m(t+H) - bid
\end{lstlisting}
This fills every order, including those the market would never have reached. The orders it wrongly
fills are exactly the ones where the price moved away before the queue was consumed---the fills a
real passive order does not get, and which would have been profitable. The naive backtest is
therefore biased toward adverse-selection-free fills that do not exist.

\subsection{Requirements of a queue-exact replay}

A replay that can answer the execution questions of this part needs:
\begin{enumerate}[leftmargin=1.5em]
\item market-by-order data, so that queue position is defined at all;
\item an order-level reconstructed book in which the simulated order occupies a place in the
  first-in-first-out queue;
\item explicit, independent latencies for outbound orders, outbound cancels and inbound market
  data;
\item an explicit fill rule stated in terms of recorded events;
\item a stated assumption about market impact.
\end{enumerate}

\begin{lstlisting}
replay:  signal at t -> order leaves at t
         -> arrives at t + feed delay + order delay
         -> appended at the tail of the real FIFO queue at its price
         -> filled only when recorded executions reach an order behind it
         -> P&L from the actual (simulated) fill time and price
\end{lstlisting}

The fill rule rests on price--time priority: if a recorded aggressor reached a resting order that
arrived later than the simulated order at the same price, it must have passed through the
simulated order first. The standard simplifying assumption is \emph{no market impact}: the
recorded market is replayed as it happened, as though the simulated order had not been there. Every
cost measured this way is a cost before the market's reaction to the order.

Replay and generative simulation answer different questions. Generative simulators produce new
market trajectories and are judged on how realistic they are; a replay simulator keeps the recorded
market fixed and produces only the counterfactual outcome of the researcher's own orders. GPU
simulators such as JAX-LOB~\citep{frey2023} move the matching engine onto accelerators to speed up
reinforcement-learning training; the concern here is not throughput but the fidelity of the fill
inference. Part~IV describes one implementation of a queue-exact replay.
